%% Requires compilation with XeLaTeX or LuaLaTeX
\documentclass[10pt,xcolor={table,dvipsnames},t]{beamer}
\usetheme{UVic}

\usepackage{booktabs}

% Reduce list spacing (beamer-native approach)
\setbeamertemplate{itemize/enumerate body begin}{\vspace{-0.5ex}}
\setbeamertemplate{itemize/enumerate subbody begin}{\vspace{-0.5ex}}

\title[LA Kelp Prototype Progress]{LA Kelp Detection Prototype}
\subtitle{Weekly Progress Report}
\author{Ekaba Bisong}
\institute{Department of Computer Science\\University of Victoria}
\date{\today}

\begin{document}

% Title Slide
\begin{frame}
  \titlepage
\end{frame}

% Table of Contents
\begin{frame}{Outline}
  \tableofcontents
\end{frame}

%% ============================================================
%% WEEK 1: Monday, February 10, 2026
%% ============================================================
\section{Monday, February 10, 2026}

\begin{frame}
\centering
\vfill
{\Large\textbf{Week 1}}\\[1em]
{\large Monday, February 10, 2026}
\vfill
\end{frame}

\begin{frame}{Week 1: Committee Meeting \& Project Setup}

\textbf{Committee Meeting (February 5, 2026):}
\begin{itemize}
  \item Received feedback to build full yet simple prototype
  \item Focus on spectral heuristics only (no ML training initially)
  \item Global Policy Learning only (L$_{\text{R-I}}$ automaton)
  \item Clarified scope: kelp \textbf{mapping}, not full ecological monitoring
\end{itemize}

\vspace{0.5em}
\textbf{Project Setup Completed:}
\begin{itemize}
  \item Created prototype directory structure
  \item Set up science log and weekly reporting system
  \item Created implementation plan
\end{itemize}

\end{frame}

\begin{frame}{Week 1: Prototype Architecture}

\textbf{Prototype Scope:}
\begin{itemize}
  \item 4 spectral heuristics: NDVI, FAI, GNDVI, Ensemble
  \item L$_{\text{R-I}}$ automaton for heuristic selection
  \item Test on existing $\sim$430 tiles from data\_bc
\end{itemize}

\vspace{0.5em}
\textbf{Infrastructure:}
\begin{itemize}
  \item config.yaml for experiment parameters
  \item Makefile for running experiments
  \item Experiment tracker with auto-logging
\end{itemize}

\end{frame}

\begin{frame}{Week 1: Next Steps}

\textbf{Immediate Tasks:}
\begin{enumerate}
  \item Implement config.yaml and Makefile
  \item Create experiment tracker
  \item Implement data loader for Sentinel-2 tiles
  \item Begin spectral heuristics implementation
\end{enumerate}

\vspace{0.5em}
\textbf{Success Criteria:}
\begin{itemize}
  \item LA correctly selects heuristics and updates probabilities
  \item Probability vector converges toward better heuristics
  \item LA outperforms random selection baseline
\end{itemize}

\end{frame}

%% ============================================================
%% WEEK 2: Monday, February 17, 2026
%% ============================================================
\section{Monday, February 17, 2026}

\begin{frame}
\centering
\vfill
{\Large\textbf{Week 2}}\\[1em]
{\large Monday, February 17, 2026}
\vfill
\end{frame}

\begin{frame}{Week 2: Prototype Implementation Complete}
\small
\textbf{All Components Implemented:}
\begin{itemize}
  \item Data loader for Sentinel-2 tiles (10 bands)
  \item 4 spectral heuristics: NDVI, FAI, GNDVI, Ensemble
  \item L$_{\text{R-I}}$ automaton with probability updates
  \item Reward function (binary/continuous modes)
  \item Experiment runner with convergence plots; Baselines: Random, Fixed, Oracle
\end{itemize}

\vspace{0.2em}
\textbf{Infrastructure:}
\begin{itemize}
  \item \texttt{make run} / \texttt{make baselines} -- run experiments
  \item Conda environment ``kelp'' (Python 3.10)
\end{itemize}

\end{frame}

\begin{frame}{Week 2: Baseline Results}
\small
\textbf{Test Set Performance (IoU):}

\vspace{0.3em}
\begin{table}
\centering
\footnotesize
\begin{tabular}{lrl}
\toprule
\textbf{Method} & \textbf{IoU} & \textbf{Role} \\
\midrule
Oracle & 0.0765 & Upper bound \\
Fixed GNDVI & 0.0699 & Best single heuristic \\
Fixed Ensemble & 0.0695 & Second best \\
Fixed NDVI & 0.0533 & Middle performer \\
Random & 0.0445 & Lower bound \\
Fixed FAI & 0.0074 & Worst heuristic \\
\bottomrule
\end{tabular}
\end{table}

\vspace{0.3em}
\textbf{Key Insight:} IoU values are low (0.007--0.076) due to class imbalance (kelp is sparse in tiles).

\end{frame}

\begin{frame}{Week 2: Experiment 1 -- Binary Reward (Failed)}
\small
\textbf{Configuration:}
\begin{itemize}
  \item \texttt{reward.type: "binary"}, \texttt{reward.iou\_threshold: 0.3}
\end{itemize}

\vspace{0.3em}
\textbf{Result:} LA converged to \textbf{FAI} (the \emph{worst} heuristic!)

\vspace{0.3em}
\textbf{Root Cause Analysis:}
\begin{itemize}
  \item Binary reward: $r = 1$ if IoU $> 0.3$, else $r = 0$
  \item Actual IoU values: 0.007 -- 0.076 (never exceed 0.3)
  \item \textcolor{red}{Almost no tiles received positive reward}
  \item Automaton received near-zero learning signal
  \item Probability updates were essentially random drift
\end{itemize}

\end{frame}

\begin{frame}{Week 2: Fix -- Continuous Reward}

\textbf{Problem:} Threshold 0.3 unreachable for this dataset.

\vspace{0.5em}
\textbf{Solution:} Use continuous reward (IoU directly as reward signal)
\begin{itemize}
  \item Changed: \texttt{reward.type: "continuous"}
  \item Now every tile provides gradient information
  \item IoU of 0.07 gives more reward than IoU of 0.007
\end{itemize}

\vspace{0.5em}
\textbf{Expected Outcome:}
\begin{itemize}
  \item LA should converge toward GNDVI/Ensemble ($\sim$0.07 IoU)
  \item LA should move away from FAI ($\sim$0.007 IoU)
  \item Test IoU should exceed random baseline (0.0445)
\end{itemize}

\end{frame}

\begin{frame}{Week 2: Experiment 2 -- Continuous Reward, $\alpha=0.1$}
\small
\textbf{Configuration:} \texttt{reward.type: "continuous"}, \texttt{automaton.alpha: 0.1}

\vspace{0.3em}
\textbf{Result:} LA converged to \textbf{Ensemble} (correct!) but \textcolor{orange}{too fast}.

\vspace{0.3em}
\begin{table}
\centering
\footnotesize
\begin{tabular}{lrr}
\toprule
\textbf{Epoch} & \textbf{Ensemble Prob} & \textbf{Entropy} \\
\midrule
1 & 98.76\% & 0.096 \\
2 & 99.9999\% & 0.000 \\
3--10 & 100\% & 0.000 \\
\bottomrule
\end{tabular}
\end{table}

\vspace{0.3em}
\textbf{Issue:} Winner-take-all dynamics. No exploration after Epoch 2.

\end{frame}

\begin{frame}{Week 2: Experiment 3 -- Lower Learning Rate, $\alpha=0.01$}
\small
\textbf{Configuration:} \texttt{reward.type: "continuous"}, \texttt{automaton.alpha: 0.01}

\vspace{0.2em}
\textbf{Result:} \textcolor{green}{Gradual convergence} -- proper learning dynamics!

\vspace{0.2em}
\begin{table}
\centering
\footnotesize
\begin{tabular}{lrrrr}
\toprule
\textbf{Epoch} & \textbf{NDVI} & \textbf{FAI} & \textbf{GNDVI} & \textbf{Ensemble} \\
\midrule
1 & 37\% & 3\% & 27\% & 34\% \\
4 & 44\% & 0.03\% & 25\% & 30\% \\
7 & 29\% & $\sim$0\% & 18\% & 53\% \\
10 & 20\% & $\sim$0\% & \textbf{41\%} & 39\% \\
\bottomrule
\end{tabular}
\end{table}

\vspace{0.2em}
\textbf{Key:} FAI eliminated, GNDVI/Ensemble competing, entropy = 1.52

\end{frame}

\begin{frame}{Week 2: Experiment Comparison}
\footnotesize
\begin{table}
\centering
\begin{tabular}{lcc}
\toprule
\textbf{Metric} & \textbf{Exp 2 ($\alpha$=0.1)} & \textbf{Exp 3 ($\alpha$=0.01)} \\
\midrule
Convergence & Instant (Epoch 2) & Gradual (10 epochs) \\
Final Entropy & 0.0 & 1.52 \\
Best Heuristic Prob & 100\% & 41\% \\
Test IoU & 0.0719 & 0.0711 \\
vs Random & +0.0151 & +0.0144 \\
Gap to Oracle & 6.0\% & 7.0\% \\
\bottomrule
\end{tabular}
\end{table}

\small
\textbf{Both experiments:}
\begin{itemize}
  \item \textcolor{green}{\checkmark} Outperform random baseline; \textcolor{green}{\checkmark} Eliminate FAI (worst)
  \item \textcolor{green}{\checkmark} Converge toward best heuristics (GNDVI/Ensemble)
\end{itemize}

\end{frame}

\begin{frame}{Week 2: Key Learnings}
\small
\textbf{Lessons from Experiments 1--3:}
\begin{enumerate}
  \item Binary reward requires threshold calibration to data distribution
  \item Continuous reward provides gradient signal for low IoU values
  \item Learning rate $\alpha$ controls convergence: $\alpha=0.1$ (instant), $\alpha=0.01$ (gradual)
  \item L$_{\text{R-I}}$ correctly learns heuristic quality ranking
\end{enumerate}

\end{frame}

\begin{frame}{Next Steps: Prototype Refinement}
\small
\textbf{Prototype Refinement:}
\begin{enumerate}
  \item Run more epochs (20--50) to see full convergence behavior
  \item Try $\alpha=0.001$ for slower/stable convergence curves
  \item Add $\varepsilon$-greedy exploration to prevent probability collapse
\end{enumerate}

\vspace{0.3em}
\textbf{LA Variant Comparison:}
\begin{enumerate}
  \item Compare L$_{\text{R-I}}$ with L$_{\text{R-P}}$ (reward-penalty)
  \item Test Pursuit algorithm (faster convergence via estimation)
  \item Test Estimator/SERI (Bayesian estimation + pursuit)
  \item Evaluate VSLA (adaptive learning rate based on entropy)
\end{enumerate}

\end{frame}

\begin{frame}{Next Steps: LA Variant Comparison}
\small
Implemented in \texttt{automaton.py}:

\vspace{0.3em}
\begin{table}
\centering
\footnotesize
\begin{tabular}{lll}
\toprule
\textbf{Algorithm} & \textbf{Description} & \textbf{Expected Behavior} \\
\midrule
L$_{\text{R-I}}$ & Reward only, ignore penalty & Conservative, stable \\
L$_{\text{R-P}}$ & Reward and penalty updates & Faster but noisier \\
Pursuit & Estimates reward, pursues best & 10--20x faster convergence \\
Discretized Pursuit & Finite probability levels & Even faster \\
Estimator (SERI) & Bayesian estimation + pursuit & Fastest $\varepsilon$-optimal \\
VSLA & Adaptive $\alpha$ based on entropy & Auto-balances explore/exploit \\
\bottomrule
\end{tabular}
\end{table}

\vspace{0.3em}
\textbf{Next:} Implement and run experiments comparing all variants on kelp detection task.

\end{frame}

%% ============================================================
%% WEEK 3: Monday, February 24, 2026
%% ============================================================
\section{Monday, March 6, 2026}

\begin{frame}
\centering
\vfill
{\Large\textbf{Week 3}}\\[1em]
{\large Monday, March 6, 2026}
\vfill
\end{frame}

\begin{frame}{Week 3: Kelp Dataset Survey}
\small
\textbf{Objective:} Identify and acquire labeled kelp datasets for LA training and validation.

\vspace{0.3em}
\textbf{Target Species:}
\begin{itemize}
  \item Giant kelp (\textit{Macrocystis pyrifera}) -- California, Baja California
  \item Bull kelp (\textit{Nereocystis luetkeana}) -- Pacific Northwest
\end{itemize}

\vspace{0.3em}
\textbf{Satellite Data Sources:}
\begin{itemize}
  \item Landsat 5, 7, 8, 9 (30 m resolution)
  \item Sentinel-2 (10 m resolution) -- higher priority for future work
\end{itemize}

\end{frame}

\begin{frame}{Week 3: Downloaded Dataset}
\footnotesize
\begin{table}
\centering
\begin{tabular}{lllr}
\toprule
\textbf{Dataset} & \textbf{Type} & \textbf{Best For} & \textbf{Size} \\
\midrule
Figshare Kelp & Polygon masks & \textbf{Training} & 73 MB \\
\bottomrule
\end{tabular}
\end{table}

\vspace{0.3em}
\textbf{Figshare Mask R-CNN Dataset:}
\begin{itemize}
  \item 432 Landsat images (train: 323, val: 30, test: 79)
  \item VIA JSON polygon annotations
  \item Performance: Dice 0.93, Jaccard 0.87
  \item Region: Southern California and Baja California
  \item Species: Giant kelp (\textit{Macrocystis pyrifera})
\end{itemize}

\vspace{0.3em}
\textbf{Note:} SBC LTER contains derived canopy products (not labeled imagery), so not useful for training/validation.

\end{frame}

\begin{frame}{Week 3: Datasets Evaluated \& Rejected}
\small
\textbf{Datasets NOT for Kelp Research:}
\begin{table}
\centering
\footnotesize
\begin{tabular}{ll}
\toprule
\textbf{Dataset} & \textbf{Actual Target (Not Kelp)} \\
\midrule
USF Floating Algae & Sargassum, Cyanobacteria, Ulva \\
Marine Debris & Plastic, driftwood, debris \\
IMO Macroalgae & European intertidal algae \\
\bottomrule
\end{tabular}
\end{table}

\vspace{0.3em}
\textbf{Key Distinction:} Kelp (Macrocystis, Nereocystis) are subtidal canopy-forming species. The above datasets target floating algae or debris.

\end{frame}

\begin{frame}{Week 3: High-Quality Datasets (Contact Required)}
\small
\textbf{1. Floating Forests (DrivenData Source)}
\begin{itemize}
  \item 5,635 labeled Landsat images; $>$1M citizen classifications
  \item Regions: Falkland Islands, California, Tasmania
  \item Dr.\ Kyle Cavanaugh (UCLA): \texttt{kcavanaugh@geog.ucla.edu}
  \item Dr.\ Jarrett Byrnes (UMass): \texttt{Jarrett.Byrnes@umb.edu}
\end{itemize}

\vspace{0.3em}
\textbf{2. Hakai Institute (BC, Canada)}
\begin{itemize}
  \item Field data + aerial imagery for BC Central Coast
  \item Species-specific labels: Macrocystis vs Nereocystis
  \item ML training pipeline: \url{github.com/HakaiInstitute/hakai-ml-train}
  \item Data catalogue: \url{catalogue.hakai.org}
\end{itemize}

\end{frame}

\begin{frame}{Week 3: Dataset Summary \& Next Steps}
\small
\textbf{Current Status:}
\begin{itemize}
  \item \textcolor{green}{\checkmark} 1 labeled dataset downloaded (Figshare)
  \item \textcolor{green}{\checkmark} Created data survey report.
\end{itemize}

\vspace{0.3em}
\textbf{Next Steps:}
\begin{enumerate}
  \item Use Figshare data for initial LA heuristic evaluation
  \item Contact Cavanaugh/Byrnes for Floating Forests labels
  \item Contact Hakai Institute for BC kelp data
  \item Follow-up with Mohsen to obtain his data and models
  \item Look for available open-source Kelp models
\end{enumerate}

\end{frame}

%% ============================================================
%% ADD NEW WEEKS BELOW (copy section template)
%% ============================================================

%% \section{Monday, March 13, 2026}
%% \begin{frame}{Week 4: [Title]}
%% ...
%% \end{frame}

\end{document}
